\chaptertitle{第五章\quad 指数对数——实数运算的完成}

\sectiontitle{第16讲\quad 从"几个2相乘"到指数增长的直觉}

\subsection*{观察一个规律}

$2^1=2$（1个2）\\
$2^2=2\times2=4$（2个2相乘）\\
$2^3=2\times2\times2=8$（3个2相乘）\\
$2^4=2\times2\times2\times2=16$（4个2相乘）\\
$2^5=2\times2\times2\times2\times2=32$（5个2相乘）

你发现了吗？每多乘一个2，结果就翻一倍。

这就是\textbf{指数增长}的雏形——不是"每次加一个固定数"，而是"每次翻倍"。

\subsection*{棋盘上的麦粒}

现在来讲那个经典的故事。

传说国际象棋的发明人向国王要赏赐。棋盘有64格，他提出的要求是：第一格放1粒麦子，第二格放2粒，第三格放4粒，以此类推，每格翻一倍。

国王一口答应了——这个要求看起来太简单了。

让我们算一算第64格需要多少粒麦子：

第一格：$1 = 2^0$\\
第二格：$2 = 2^1$\\
第三格：$4 = 2^2$\\
第四格：$8 = 2^3$\\
$\vdots$\\
第64格：$2^{63}$

$2^{63}$ 是多少？大约 9,220,000,000,000,000,000（92亿亿）。全世界所有的麦田加起来都不够。

这就是指数增长的力量：\textbf{开始慢到你看不出来，后面快到你想不到}。

\begin{example}
$2^6$ 是 $2^5$ 的几倍？
\end{example}
\begin{solution}
$2^5=32$，$2^6=64$。$64\div32=2$，所以刚好是2倍。每次指数加1，值翻一倍。
\end{solution}

\begin{example}
从 $2^1$ 开始，写到 $2^{10}$，看看2的幂的个位有什么规律。
\end{example}
\begin{solution}
$2^1=2$，$2^2=4$，$2^3=8$，$2^4=16$，$2^5=32$，$2^6=64$，$2^7=128$，$2^8=256$，$2^9=512$，$2^{10}=1024$。
个位依次是：2,4,8,6,2,4,8,6,2,4——每4个一循环！
\end{solution}

\begin{example}
$3^4$ 比 $3^3$ 大多少？
\end{example}
\begin{solution}
$3^3=27$，$3^4=81$。$81-27=54$，相差54倍。不对——应该是 $81\div27=3$，是3倍的关系。指数增加1，比值固定的（这里是乘以3），但差值不是固定的。
\end{solution}

\subsection*{10的幂的规律}

$10^1=10$\\
$10^2=100$\\
$10^3=1000$\\
$10^4=10000$

规律很简单：10的n次方，就是1后面跟n个0。

那反过来呢？$10^0=1$（1后面跟0个0）。$10^{-1}=0.1$。$10^{-2}=0.01$。

指数可以是0，也可以是负数——这是指数概念向两端的扩展。

\begin{exercise}
\item 计算下列乘方，观察规律：
  \begin{subexercise}
    \item $2^0,\;2^1,\;2^2,\;2^3,\;2^4,\;2^5,\;2^6$
    \item 每次后一个数是前一个数的几倍？
  \end{subexercise}

\item 计算：
  \begin{subexercise}
    \item $2^6$
    \item $2^7$
    \item $2^8$
    \item $2^9$
    \item $2^{10}$
  \end{subexercise}

\item 10的幂：
  \begin{subexercise}
    \item $10^1$
    \item $10^3$
    \item $10^5$
    \item $10^0$
    \item $10^{-1}$
    \item $10^{-2}$
  \end{subexercise}

\item 一张纸对折1次变成2层，对折2次变成4层$\ldots$ 对折多少次后层数超过1000？
  （提示：$2^{10}=1024$）

\item 假如你今天存1元，第二天2元，第三天4元，每天翻一倍。
  \begin{subexercise}
    \item 第10天你要存多少元？
    \item 第20天呢？
    \item 第30天呢？（$2^{30}\approx10.7亿$）
  \end{subexercise}

\item 观察 $10^3=1000$，$10^2=100$，$10^1=10$，$10^0=1$，$10^{-1}=0.1$。\\
  指数和结果的小数点位置有什么关系？
\end{exercise}

\sectiontitle{第17讲\quad 对数的登场——"几次方"的问题}

\subsection*{把问题反过来}

我们知道 $2^3=8$。现在反过来问：\textbf{2的几次方等于8？}

答案是3。因为 $2^3=8$。

这个问题在数学上有一个专门的符号：

\[
\log_2 8 = 3
\]

读作"以2为底8的对数等于3"。\textbf{对数的本质就是问"几次方"}。

每一次乘方 $a^b=c$ 都有三个角色：底数 $a$、指数 $b$、幂 $c$。
\begin{itemize}
  \item 开方问的是：已知幂和指数，求底数（$\sqrt[b]{c}=a$）
  \item 对数问的是：已知幂和底数，求指数（$\log_a c = b$）
\end{itemize}

\begin{example}
$\log_2 4$ 等于多少？
\end{example}
\begin{solution}
2的几次方等于4？$2^2=4$，所以 $\log_2 4 = 2$。
\end{solution}

\begin{example}
$\log_3 81$ 等于多少？
\end{example}
\begin{solution}
3的几次方等于81？$3^4=81$，所以 $\log_3 81 = 4$。
\end{solution}

\subsection*{对数是"找指数"的运算}

你也许觉得对数符号奇怪，但它做的事情你已经会了——只是换个问法：

"2的几次方等于8？" → "$\log_2 8$"

"3的几次方等于9？" → "$\log_3 9$"

"10的几次方等于1000？" → "$\log_{10} 1000$"

每道题你都可以先在心里想"$?$ 的几次方"，然后把问号写在 $\log$ 的右下角。

\begin{example}
$\log_5 25$ 等于多少？
\end{example}
\begin{solution}
在心里想：5的几次方等于25？5的2次方。所以 $\log_5 25 = 2$。
\end{solution}

\begin{example}
$\log_2 1$ 等于多少？
\end{example}
\begin{solution}
2的几次方等于1？任何非零数的0次方等于1，$2^0=1$。所以 $\log_2 1 = 0$。
\end{solution}

\begin{exercise}
\item 口算（想"几的几次方"）：
  \begin{subexercise}
    \item $\log_2 4$
    \item $\log_2 8$
    \item $\log_2 16$
    \item $\log_2 32$
    \item $\log_2 64$
    \item $\log_2 1$
    \item $\log_2 2$
    \item $\log_2 128$
  \end{subexercise}

\item 口算：
  \begin{subexercise}
    \item $\log_3 9$
    \item $\log_3 27$
    \item $\log_3 81$
    \item $\log_5 25$
    \item $\log_4 16$
    \item $\log_4 64$
    \item $\log_7 7$
    \item $\log_8 1$
  \end{subexercise}

\item 把指数式改写成对数式：
  \begin{subexercise}
    \item $2^5=32$
    \item $3^3=27$
    \item $10^4=10000$
    \item $5^2=25$
    \item $2^{10}=1024$
    \item $4^3=64$
  \end{subexercise}

\item 把对数式改写成指数式：
  \begin{subexercise}
    \item $\log_2 64 = 6$
    \item $\log_3 81 = 4$
    \item $\log_{10} 100 = 2$
    \item $\log_5 125 = 3$
    \item $\log_7 49 = 2$
    \item $\log_4 256 = 4$
  \end{subexercise}

\item 填空并说说规律：
  \begin{subexercise}
    \item $\log_a 1 = \underline{\qquad}$（任何 $a>0$ 且 $a\neq1$）
    \item $\log_a a = \underline{\qquad}$
  \end{subexercise}

\item 思考：$\log_2 3$ 大概是多少？
  （提示：$2^1=2$，$2^2=4$，所以 $\log_2 3$ 在1和2之间。）
\end{exercise}

\sectiontitle{第18讲\quad 常用对数——以10为底的对数}

\subsection*{为什么以10为底特别重要？}

我们用的数字是十进制。所以在很多实际问题中，以10为底的对数最方便。它有一个专属简写：$\lg$。

\[
\log_{10} x = \lg x
\]

\begin{example}
$\lg 10 = 1$（因为 $10^1=10$）\\
$\lg 100 = 2$（因为 $10^2=100$）\\
$\lg 1000 = 3$（因为 $10^3=1000$）\\
$\lg 1 = 0$（因为 $10^0=1$）\\
$\lg 0.1 = -1$（因为 $10^{-1}=0.1$）
\end{example}

规律：$\lg$ 告诉你的就是这个数有几位。

\subsection*{生活中的对数}

你可能已经用过对数，只是不自知：

\textbf{声音——分贝（dB）：}
\begin{itemize}
  \item 耳语约 20dB → 能量是基准的 $10^2=100$ 倍
  \item 正常交谈约 60dB → 能量是基准的 $10^6=100万$ 倍
  \item 摇滚演唱会约 120dB → 能量是基准的 $10^{12}=1万亿$ 倍
\end{itemize}

把巨大的能量范围（1到1万亿），压缩成了好读的小数字（0到120）——这就是对数的力量。

\textbf{地震——里氏震级：}
\begin{itemize}
  \item 5级地震的能量大约是4级的10倍
  \item 5级和7级差多少？$10^{7-5}=100$倍
\end{itemize}

\textbf{酸碱度——pH值：}
\[
pH = -\lg[H^+]
\]
pH=7是中性，pH=6的酸性是pH=7的10倍。

\begin{example}
一个声音的强度是基准的 $10^8$ 倍，它是多少分贝？
\end{example}
\begin{solution}
$10^8$ 倍 → $\lg(10^8) = 8$，但分贝的定义要乘以10，所以是 80dB。
\end{solution}

\begin{example}
$\lg 5000$ 大概在多少到多少之间？
\end{example}
\begin{solution}
$1000 < 5000 < 10000$，所以 $\lg 1000 < \lg 5000 < \lg 10000$。
$\lg 1000 = 3$，$\lg 10000 = 4$，所以 $\lg 5000$ 在3到4之间。
\end{solution}

\begin{example}
如果 $\lg x = 2.5$，那么 $x$ 大概是多少？
\end{example}
\begin{solution}
$\lg x = 2.5$ 的意思是 $10^{2.5} = x$。$10^2 = 100$，$10^3 = 1000$，所以 $x$ 应该在100到1000之间。
$10^{2.5} = 10^{2} \times 10^{0.5} \approx 100 \times 3.16 = 316$。
\end{solution}

\begin{example}
一个6级地震的能量是4级地震的多少倍？
\end{example}
\begin{solution}
每升1级能量变为10倍。从4级到6级升了2级，所以是 $10^2 = 100$倍。
\end{solution}

\begin{exercise}
\item 口算 $\lg$：
  \begin{subexercise}
    \item $\lg 10$
    \item $\lg 100$
    \item $\lg 10000$
    \item $\lg 100000$
    \item $\lg 1$
    \item $\lg 0.1$
    \item $\lg 0.01$
    \item $\lg 0.001$
  \end{subexercise}

\item 填空：
  \begin{subexercise}
    \item $\lg 10^3 = \underline{\qquad}$
    \item $\lg 10^5 = \underline{\qquad}$
    \item $\lg 10^{-2} = \underline{\qquad}$
    \item $\lg 10^n = \underline{\qquad}$
  \end{subexercise}

\item 一个声音的强度是基准的 $10^5$ 倍，有多少分贝？

\item 地震震级每升1级，能量变为原来的10倍。
  \begin{subexercise}
    \item 4级地震的能量是3级的多少倍？
    \item 6级地震的能量是3级的多少倍？
    \item 8级地震的能量是5级的多少倍？
  \end{subexercise}

\item 已知 $\lg 2 \approx 0.3010$，$\lg 3 \approx 0.4771$。
  \begin{subexercise}
    \item 你猜 $\lg 6$ 大概是多少？（提示：$6=2\times3$）
    \item $\lg 5$ 呢？（提示：$5=10\div2$）
  \end{subexercise}

\item 判断题：$\lg$ 的结果如果是整数，那么原数一定是10的整数次幂。对吗？
\end{exercise}

\sectiontitle{第19讲\quad 对数的魔法——化乘为加}

\subsection*{做一个发现}

先算几道简单的：

$\lg 10 = 1$，$\lg 100 = 2$，$\lg(10\times100) = \lg 1000 = 3$。

现在看：$\lg 10 + \lg 100 = 1 + 2 = 3$。

$\lg(10\times100)$ 和 $\lg 10 + \lg 100$ 相等！

再来一个：$\lg(100\times1000) = \lg 100000 = 5$；$\lg 100 + \lg 1000 = 2 + 3 = 5$。

\textbf{这不是巧合！}

\subsection*{对数的核心性质}

\[
\log_a (M\times N) = \log_a M + \log_a N
\]

\[
\log_a \left(\frac{M}{N}\right) = \log_a M - \log_a N
\]

**对数把乘法变成加法，把除法变成减法。**

如果你觉得这没什么了不起，想象一下：在500年前没有计算器的时代，天文学家要计算两个很大的数相乘，需要算整整一天。但有了对数表，他们只需要查表找到这两个数的对数，做一次加法，再查表把结果变回去——几分钟就完成了。

法国数学家拉普拉斯说过：**"对数的发明，延长了天文学家的寿命。"**

\begin{example}
已知 $\lg 2\approx0.3010$，$\lg 3\approx0.4771$，求 $\lg 6$ 和 $\lg 1.5$。
\end{example}
\begin{solution}
$\lg 6 = \lg(2\times3) = \lg 2 + \lg 3 \approx 0.3010+0.4771 = 0.7781$。\\
$\lg 1.5 = \lg(3\div2) = \lg 3 - \lg 2 \approx 0.4771-0.3010 = 0.1761$。
\end{solution}

\begin{example}
计算：$\lg 5 + \lg 2$
\end{example}
\begin{solution}
$\lg 5 + \lg 2 = \lg(5\times2) = \lg 10 = 1$。
\end{solution}

\subsection*{指数和对数——互逆运算的最后一块拼图}

现在我们可以回头看看乘方的两个逆运算了：

\[
a^b = c \quad\Longrightarrow\quad
\begin{cases}
\sqrt[b]{c} = a & \text{开方——已知幂和指数，求底数}\\
\log_a c = b & \text{对数——已知幂和底数，求指数}
\end{cases}
\]

开方和对数，一个是乘方的第一个逆运算，一个是第二个逆运算。两者合在一起，才构成了乘方的完整逆运算。

\begin{exercise}
\item 用对数的性质计算：
  \begin{subexercise}
    \item $\lg 2 + \lg 5$
    \item $\lg 25 + \lg 4$
    \item $\lg 1000 - \lg 100$
    \item $\lg 0.1 + \lg 100$
    \item $\lg 50 + \lg 2$
    \item $\lg 8 - \lg 2$
  \end{subexercise}

\item 已知 $\lg 2 \approx 0.3010$，$\lg 3 \approx 0.4771$，求：
  \begin{subexercise}
    \item $\lg 12$（提示：$12=3\times4=3\times2^2$）
    \item $\lg 18$
    \item $\lg 30$
    \item $\lg 5$
    \item $\lg 1.5$
    \item $\lg 0.5$
  \end{subexercise}

\item 为什么 $\lg 25 + \lg 4 = 2$？写出推理过程。

\item 用文字解释：为什么对数能把乘法变成加法？

\item 如果 $\log_a M = 3$，$\log_a N = 4$，那么 $\log_a (MN) = \underline{\qquad}$。

\item 思考题：$\log_a a^n = \underline{\qquad}$？（用一句话回答）

\item 查表题（感受一下古代天文学家的工具）：\\
  下表是 $\lg$ 的简表：
  \[
  \begin{array}{c|c}
  x & \lg x \\ \hline
  1 & 0.0000 \\
  2 & 0.3010 \\
  3 & 0.4771 \\
  4 & 0.6021 \\
  5 & 0.6990 \\
  6 & 0.7782 \\
  7 & 0.8451 \\
  8 & 0.9031 \\
  9 & 0.9542 \\
  10 & 1.0000
  \end{array}
  \]
  用这张表计算 $2\times3$：
  \begin{enumerate}
    \item 查 $\lg 2$ 和 $\lg 3$，相加
    \item 和是 0.7781，在表中对应哪个数？
    \item $2\times3$ 等于多少？
  \end{enumerate}
  这就是对数的原始用法。
\end{exercise}
\end{parameter>
